% Generated draft: requires mathematical and source review.
\documentclass[11pt]{article}
\usepackage[margin=25mm]{geometry}
\usepackage{fontspec}
\setmainfont{Latin Modern Roman}
\usepackage{amsmath,amssymb,mathtools,unicode-math}
\setmathfont{Latin Modern Math}
\usepackage{enumitem,needspace,xurl,hyperref}
\setlength{\parindent}{0pt}
\setlength{\parskip}{5pt}
\setlength{\emergencystretch}{3em}
\newcommand{\N}{\mathbb{N}}
\newcommand{\Z}{\mathbb{Z}}
\newcommand{\Q}{\mathbb{Q}}
\newcommand{\R}{\mathbb{R}}
\newcommand{\C}{\mathbb{C}}
\newcommand{\F}{\mathbb{F}}
\newcommand{\A}{\mathbb{A}}
\newcommand{\PP}{\mathbb{P}}
\newcommand{\E}{\mathbb{E}}
\newcommand{\eps}{\varepsilon}
\newcommand{\dd}{\,\mathrm{d}}
\begin{document}
\section*{UnsolvedMath ID: 800617}
\textbf{Source problem code:} catalogue entry
\subsection{Definitions and mathematical statement}
Let \(\mathfrak c=|\R|=2^{\aleph_0}\). The Continuum Hypothesis (CH) is the assertion \(\mathfrak c=\aleph_1\). Regard \(\Z^{\N}\) as an additive group with coordinatewise addition. For the topological formulation below, equip \(\Z\) with the discrete topology, \(\Z^{\N}\) with the product topology, and every subgroup with the subspace topology. This is the standard intended topological interpretation; the supplied excerpt does not explicitly specify the topology.

For a topological space \(X\), an \(\omega\)-cover is a family \(\mathcal U\) of subsets of \(X\) such that \(X\notin\mathcal U\) and every finite subset of \(X\) is contained in some member of \(\mathcal U\). A \(\gamma\)-cover is an infinite family \(\mathcal U\) of subsets of \(X\) such that every point of \(X\) belongs to all but finitely many members of \(\mathcal U\). Using countable-cover conventions, let \(\Omega(X)\) and \(\Gamma(X)\) denote, respectively, the collections of countable open \(\omega\)-covers and countably infinite open \(\gamma\)-covers, with \(X\) excluded as a cover member. Let \(\mathrm B_\Omega(X)\) and \(\mathrm B_\Gamma(X)\) be the corresponding collections with Borel subsets of \(X\) in place of open subsets. Here Borel always refers to the topology of \(X\). The excerpt does not supply the source's definitions of these cover classes; these standard conventions require confirmation against the cited source.

For collections \(\mathcal A\) and \(\mathcal B\) of covers of \(X\), the selection principle \(\mathsf S_1(\mathcal A,\mathcal B)\) means that for every sequence \((\mathcal U_n)_{n\in\N}\) with \(\mathcal U_n\in\mathcal A\), there are choices \(U_n\in\mathcal U_n\) such that \(\{U_n:n\in\N\}\in\mathcal B\). In particular, when \(\mathcal B\) is a class of \(\gamma\)-covers, the selected family must be infinite.

Under these conventions, the question is whether CH implies that there exists a subgroup \(G\leq\Z^{\N}\) such that
\[
 |G|=\mathfrak c
 \quad\text{and}\quad
 \mathsf S_1\bigl(\mathrm B_\Omega(G),\mathrm B_\Gamma(G)\bigr).
\]
The source also asks for the weaker target: does CH at least imply that there exists a subgroup \(G\leq\Z^{\N}\) with
\[
 |G|=\mathfrak c
 \quad\text{and}\quad
 \mathsf S_1\bigl(\Omega(G),\Gamma(G)\bigr)?
\]
The weaker question is a separate existence target; it does not prescribe a particular subgroup. Both questions are understood over the usual background theory ZFC, which is not explicitly stated in the supplied excerpt.
\subsection{Short English statement}
Assuming the Continuum Hypothesis, must there be an additive subgroup of \(\Z^{\N}\) of size continuum satisfying the one-element selection principle from Borel \(\omega\)-covers to Borel \(\gamma\)-covers? The source also asks whether such a subgroup must at least exist for the corresponding open-cover principle. The topology and cover conventions used here require confirmation from the original source.
\subsection{Sources}
\begin{itemize}
\item Local UnsolvedMath registry, permanent ID \(800617\), source problem code \texttt{catalogue entry}; supplied title: \emph{Selection Principles in Mathematics --- Problem 45 ([o-bdd])}. The registry records Tsaban as proposer and \(2006\) as the proposed year.
\item Tsaban, \emph{Selection Principles and Special Sets of Reals: Open Problems} (2006), Problem 45. Supplied source URL: \url{https://arxiv.org/abs/math/0609601}. The provenance records access on 2026-10-16 and extraction by \texttt{arxiv\_latex}; these are supplied provenance details, not a claim of access or verification in this draft.
\item The supplied attribution/context marker is \texttt{[o-bdd]}; its expansion and bibliographic role are not provided. The registry labels the problem \texttt{open}; this is a supplied status label, not an independently verified literature assessment.
\end{itemize}
\end{document}
